\documentclass[11pt, a4paper]{article}

\usepackage[T1]{fontenc}
\usepackage[utf8]{inputenc}
\usepackage{tgpagella}
\usepackage{tgcursor}
\usepackage{microtype}
\usepackage[spanish, es-noshorthands]{babel}
\usepackage[margin=2.5cm]{geometry}
\usepackage{fancyhdr}
\usepackage[most]{tcolorbox}
\usepackage{listings}

\lstset{
    language=Python,
    basicstyle=\ttfamily\small,
    backgroundcolor=\color{gray!5},
    frame=single,
    rulecolor=\color{gray!40},
    keywordstyle=\color{blue}\bfseries,
    stringstyle=\color{red!70!black},
    commentstyle=\color{green!60!black},
    showstringspaces=false,
    breaklines=true,
    breakatwhitespace=true
}

\pagestyle{fancy}
\fancyhf{}
\setlength{\headheight}{16pt}
\lhead{\textbf{IMT342}}
\chead{Código Opcional: Rango Jacobiano}
\rhead{Semana 8 - Sesión 2}

\definecolor{ucbblue}{RGB}{0, 51, 102}

\begin{document}

\begin{center}
    \Large\textbf{\textcolor{ucbblue}{Verificación Computacional: Rango y Singularidad 2R}}\\[0.2cm]
\end{center}

Este script refuerza conceptualmente la relación entre la geometría del robot, la evaluación numérica del Jacobiano y el cálculo algorítmico del rango matricial y su determinante[cite: 12].

\begin{lstlisting}
import numpy as np

def jacobian_2r(q1, q2, a1=1.0, a2=1.0):
    s1 = np.sin(q1)
    c1 = np.cos(q1)
    s12 = np.sin(q1 + q2)
    c12 = np.cos(q1 + q2)

    return np.array([
        [-a1*s1 - a2*s12, -a2*s12],
        [ a1*c1 + a2*c12,  a2*c12]
    ], dtype=float)

q_regular = (0.0, np.pi/2)
q_singular = (0.0, 0.0)

J_regular = jacobian_2r(*q_regular)
J_singular = jacobian_2r(*q_singular)

for name, J in [("regular", J_regular), ("singular", J_singular)]:
    print(f"\n{name}")
    print(np.round(J, 4))
    print("rank =", np.linalg.matrix_rank(J))
    print("det =", np.linalg.det(J))
\end{lstlisting}

\begin{tcolorbox}[colback=white, colframe=ucbblue, title=\textbf{Interpretación Física de la Salida[cite: 12]}]
\textbf{Caso Regular:} Muestra Rango = 2 y Determinante = 1. Posee dos direcciones instantáneas de velocidad planar independientes. \\
\textbf{Caso Singular:} Muestra Rango = 1 y Determinante = 0. Ha perdido una dirección de velocidad instantánea independiente en el espacio de tareas. El uso del determinante aquí es matemáticamente válido únicamente porque el sub-Jacobiano de traslación extraído es cuadrado ($2 \times 2$).
\end{tcolorbox}

\end{document}
